
\medskip

Pour chaque affirmation répondre par vrai ou faux. Justifier chaque réponse.

\medskip

\textbf{Affirmation 1 :} 50\,\% de \np{10350} c'est \np{10300}.

On a  $\dfrac{50}{100} \times \np{10350} = \dfrac{\np{10350}}{2} = \np{5175}$ : affirmation fausse.

\medskip

\textbf{Affirmation 2 :} $\dfrac73$ est la forme irréductible de $\dfrac{42}{18}$.

$\dfrac{42}{18} = \dfrac{7 \times 6}{3 \times 6} = \dfrac73$ : affirmation vraie.
\medskip

\textbf{Affirmation 3 :} L'équation $2x - 4 = -x + 5$ a pour solution $3$.

$2x - 4 = -x + 5$ donne en ajoutant à chaque membre $x + 4$ : $3x = 9$, soit $3 \times x = 3 \times 3$, d'où $x = 3$ : affirmation vraie.

\medskip

\textbf{Affirmation 4 :} L'arrondi à l'unité près du volume d'une boule de diamètre $21,6$~cm est \np{42213}~cm$^3$.

%\smallskip
%
%\emph{On rappelle la formule du volume d'une boule } $V = \dfrac43 \pi R^3$.
On a $V = \dfrac43 \pi \times 10,8^3 = \dfrac43 \pi \times 10,8^2 \times 10,8 = 4 \times 10,8^2 \times 3,6 \pi = \np{1679,62}\pi$. La calculatrice donne $V \approx \np{5276,7}$ soit \np{5277} cm$^3$ à l'unité près : affirmation fausse
\medskip

\begin{minipage}{0.52\linewidth}
\textbf{Affirmation 5 :} Dans la figure codée ci-contre, la  mesure de l'angle $\widehat{\text{DNB}}$, arrondie à l'unité près, est $18\degres{}$.
\end{minipage}\hfill
\begin{minipage}{0.38\linewidth}
\psset{unit=1cm}
\begin{pspicture}(5.5,2.4)
\pspolygon(0.4,0.2)(5.1,0.2)(0.4,1.9)%BND
\psframe(0.4,0.2)(0.6,0.4)
\uput[dl](0.4,0.2){B} \uput[dr](5.1,0.2){N} \uput[u](0.4,1.9){D} \uput[l](0.4,1.05){4 cm} \uput[d](2.75,0.2){12 cm}
\end{pspicture}
\end{minipage}

On a $\tan \widehat{\text{DNB}} = \dfrac{\text{DB}}{\text{NB}} = \dfrac{4}{12} = \dfrac13$ : la calculatrice donne $\widehat{\text{DNB}} \approx 18,4$ : affirmation vraie.
\bigskip

\textbf{Affirmation 6 :} On peut composer 6 codes différents avec un cadenas à 3 chiffres qui respecte les conditions suivantes :

\begin{itemize}
\item[$\bullet~~$]les deux premiers chiffres sont choisis parmi 1 ; 2 et 3 ;
\item[$\bullet~~$]un chiffre peut apparaître deux fois;
\item[$\bullet~~$]le dernier chiffre est 6.
\end{itemize}

En premier chiffre on a 3 choix et en deuxième 3 choix aussi, soit $3 \times 3 = 9$ codes différents : affirmation fausse.
\bigskip

